% ===================================================================
%  TABEL Nr. 12 — Het station. — De spoorwegen. — De schepen.
%  Traduction 2026 d'apres texto/io, controlee sur texto/fr et
%  texto/es. Consigne : voir texto/nl/10-tabel-01.tex.
% ===================================================================

%%K t12-apar-1 apar
\VUsaut{4.0mm}
\VUtitre{15.0pt}{60}{TABEL Nr. 12}
\VUsaut{2.4mm}
\VUpk{11.4pt}{40}{Het station. --- De spoorwegen. --- De schepen.}
\VUsaut{3.4mm}
\textsc{Nota.} --- \textit{De dienstregelingen die in elk land gebruikt worden zijn
verschillend, wij gebruiken dus als voorbeeld de uren van de}
\VUgras{Franse tabel.}

%%K t12-01-1 p
1. --- Ik heb zojuist door Duitsland gereisd. Voor het vertrek kocht ik een
\VUgras{spoorboekje} \textsuperscript{(15)} om het uur van vertrek van de treinen
te weten. Nadat ik vastgesteld had dat de gunstigste trein om negen uur 's avonds uit Parijs
zou vertrekken en om één uur dertien in Straatsburg zou aankomen, bereidde ik mijn reis voor,
en de volgende dag liet ik mij naar het station brengen.

%%K t12-02-1 p
2. --- Toen ik de grote vertrekhal binnenging, achter een oude plattelands\VUgras{priester}
\textsuperscript{(2)}, die zijn \VUgras{reistas} \textsuperscript{(3)} in de hand droeg, waren
al veel \VUgras{reizigers} \textsuperscript{(1)} aangekomen.

%%K t12-02-2 p
Omdat ik een \VUgras{telegram} \textsuperscript{(5)} wilde versturen, liet ik mij door een
\VUgras{controleur} \textsuperscript{(6)} \VUgras{het telegraafkantoor} \textsuperscript{(4)}
wijzen.

%%K t12-03-1 p
3. --- Voordat ik de \VUgras{wachtkamer} \textsuperscript{(7)} binnenging ging ik mij van een
\VUgras{kaartje} \textsuperscript{(9)} heen en terug voorzien.

%%K t12-04-1 p
4. --- Ik wachtte niet lang aan het \VUgras{loket} \textsuperscript{(8)}, want er waren weinig
mensen. Een \VUgras{soldaat} \textsuperscript{(10)}, die met verlof vertrok, was zo beleefd
mij zijn beurt af te staan, zodat ik genoeg tijd had om enkele inlichtingen aan de
\VUgras{stationschef} \textsuperscript{(13)} te vragen, die met de \VUgras{commissaris}
\textsuperscript{(12)} van het toezicht en met de dienstdoende \VUgras{gendarme}
\textsuperscript{(11)} sprak.

%%K t12-tit-1 sub
\VUpk{10.2pt}{40}{Het inschrijven van de bagage.}
\VUsaut{3.0mm}

%%K t12-05-1 p
5. --- Nadat ik een krant bij de \VUgras{boekhandel} \textsuperscript{(14)} gekocht had, liet
ik mijn \VUgras{bagage} \textsuperscript{(17)} wegen die een
\VUgras{kruier} \textsuperscript{(16)} zojuist op een \VUgras{bagagekar}
\textsuperscript{(19)} aangevoerd had; \VUgras{de beambte} \textsuperscript{(22)} die met de
\VUgras{weegbrug} \textsuperscript{(21)} belast is deelde mij mee, dat mijn koffer
vijfendertig kilogram weegt; dat maakte een
\VUgras{overwicht} van vijf kilogram; waarvoor ik een bijbetaling verschuldigd
ben, aan het \VUgras{loket} van de bagage \textsuperscript{(23)}.

%%K t12-06-1 p
6. --- Omdat ik te vroeg was, had ik vrije tijd om het verkeer van de reizigers in het station
te bekijken. Sommigen gaven hun handbagage af of haalden die af bij het \VUgras{bagagedepot}
\textsuperscript{(25)}; anderen raadpleegden de
\VUgras{dienstregelingen} \textsuperscript{(26)}. De aankomende reizigers
haastten zich naar de \VUgras{uitgang} \textsuperscript{(32)} met hun
\VUgras{valies} \textsuperscript{(24)} in de hand. Sommigen wachtten bij de
\VUgras{bagage-uitgifte} \textsuperscript{(27)} en toonden hun \VUgras{bewijs}
\textsuperscript{(29)} om hun koffers, \VUgras{kisten} \textsuperscript{(28)} of
\VUgras{manden} \textsuperscript{(30)} terug te krijgen.

%%K t12-07-1 p
7. --- \VUgras{De beambten van de accijns} \textsuperscript{(33)} onderzochten zorgvuldig de
pakken om vast te stellen of men iets aan te geven heeft.

%%K t12-08-1 p
8. --- Sommige reizigers gingen naar het \VUgras{buffet} \textsuperscript{(34)}, waar een
\VUgras{kelner} \textsuperscript{(35)} zich beijverde hen te bedienen. Zij moeten zich
haasten, want de \VUgras{trein} \textsuperscript{(36)}, die een sneltrein is, zal niet lang in
het station blijven.

%%K t12-09-1 p
9. --- Daarna ging ik naar het vertrekperron en zocht een doorgaande
\VUgras{wagen} \textsuperscript{(37)} om niet gedwongen te zijn tijdens de reis
van wagen te wisselen. Ik stapte in een corridorrijtuig dat een
\VUgras{coupé} \textsuperscript{(38)} voor rokers bevat en een coupé die voor
« dames alleen » gereserveerd is.

%%K t12-tit-2 sub
\VUpk{10.2pt}{40}{Vertrek in de nacht.}
\VUsaut{3.0mm}

%%K t12-10-1 p
10. --- Nadat ik de \VUgras{treeplank} \textsuperscript{(40)} beklommen, het
\VUgras{portier} \textsuperscript{(39)} gesloten, het \VUgras{gordijntje}
\textsuperscript{(41)} opgerold en mijn handbagage op het \VUgras{net} \textsuperscript{(43)}
gelegd had, ging ik op de \VUgras{bank} \textsuperscript{(42)} zitten. Toen ik de
\VUgras{lampenist} \textsuperscript{(44)} de lampen zag aansteken, dacht ik, dat ik een nacht
in de wagen zou moeten doorbrengen en ik riep de beambte die met de verhuur van de
\VUgras{hoofdkussens} \textsuperscript{(45)} belast is om er een te vragen.

%%K t12-11-1 p
11. --- De conducteur van de trein kwam de kaartjes controleren. Ik vroeg hem waarom wij nog
niet vertrokken waren, want het is het juiste uur. Hij antwoordde mij dat de
\VUgras{treinchef} \textsuperscript{(46)} bezig is fietsen in de \VUgras{bagagewagen}
\textsuperscript{(47)} te laten zetten. Bovendien had men nog niet alle voetwarmers
\textsuperscript{(48)} verwisseld.

%%K t12-12-1 p
12. --- Maar wij zouden weldra vertrekken, want de \VUgras{stoker} \textsuperscript{(50)}, op
zijn \VUgras{tender} \textsuperscript{(49)} nam kolen om het vuur van de \VUgras{locomotief}
\textsuperscript{(54)} aan te wakkeren, en
\VUgras{de machinist} \textsuperscript{(51)} wachtte alleen op het sein van de
\VUgras{onderstationschef} \textsuperscript{(52)}, die op het \VUgras{perron}
\textsuperscript{(53)} stond.

%%K t12-13-1 p
13. --- De \VUgras{ketel} \textsuperscript{(56)} van de locomotief gromde dof; de
\VUgras{schoorsteen} \textsuperscript{(55)} braakte stromen rook uit vermengd met
\VUgras{stoom} \textsuperscript{(60)}. Een schril \VUgras{gefluit}
\textsuperscript{(59)} weerklonk en de \VUgras{zuigers} \textsuperscript{(58)} begonnen te
bewegen, en dreven de \VUgras{wielen} \textsuperscript{(57)} van de zware machine aan.

%%K t12-tit-3 sub
\VUsaut{3.0mm}
\VUpk{10.2pt}{40}{Vertrek!}
\VUsaut{3.0mm}

%%K t12-14-1 p
14. --- De snelheid van de trein, die over de \VUgras{rails} \textsuperscript{(64)} liep werd
versneld; de \VUgras{perrons} \textsuperscript{(65)} verdwenen; wij passeerden de
\VUgras{privaten} \textsuperscript{(66)}, en ook een rij wagens die op het andere
\VUgras{spoor} \textsuperscript{(63)} gerangeerd stonden: \VUgras{slaapwagens}
\textsuperscript{(67)}, een \VUgras{restauratiewagen} \textsuperscript{(68)}, een
\VUgras{postwagen} \textsuperscript{(69)}.

%%K t12-noto-1 noto
\VUnotes{143.2mm}{%
\textit{(*) Nota.} --- \textit{Vanzelfsprekend is de beschrijving van de reis volgens de grens
van voor de oorlog.}%
}

%%K t12-15-1 p
15. --- Daarna trok het \VUgras{goederenstation} \textsuperscript{(71)} voor onze ogen
voorbij. Onder de loodsen laadden beambten de \VUgras{goederen} \textsuperscript{(72)} die met
een goederentrein verzonden worden.
\VUgras{Rangeerders} \textsuperscript{(76)} bij het \VUgras{waterreservoir}
\textsuperscript{(73)} rangeerden \VUgras{veewagens} \textsuperscript{(75)} en
\VUgras{platte wagens} \textsuperscript{(79)} om een \VUgras{goederentrein}
\textsuperscript{(80)} samen te stellen.

%%K t12-16-1 p
16. --- Wij passeerden de laatste \VUgras{wissel} \textsuperscript{(78)}, waarbij de
wisselwachter stond; wij kwamen langs de \VUgras{seinschijf} \textsuperscript{(86)} en het
station verdween in de verte.

%%K t12-17-1 p
17. --- Weldra reden wij over een \VUgras{ijzeren brug} \textsuperscript{(74)} die over de
\VUgras{stroom} \textsuperscript{(81)} loopt. Nadat wij die brug overgestoken waren, volgden
wij de rivier, waarop krachtige
\VUgras{sleepboten} \textsuperscript{(82)} zware \VUgras{vrachtschuiten}
\textsuperscript{(83)} trokken. Wij zagen bovendien een \VUgras{passagiersschip}
\textsuperscript{(84)}, toen het bij een \VUgras{ponton} \textsuperscript{(85)} aanlegde.

%%K t12-18-1 p
18. --- Nadat wij razendsnel verscheidene stations gepasseerd waren, reden wij door enkele
\VUgras{tunnels} \textsuperscript{(87)} en lieten achter ons ontelbare \VUgras{huisjes}
\textsuperscript{(88)} van \VUgras{overwegwachters}. Gewiegd door het schokken van de trein
viel ik in een diepe slaap.

%%K t12-tit-4 sub
\VUsaut{3.0mm}
\VUpk{10.2pt}{40}{Station Deutsch-Avricourt.}
\VUsaut{3.0mm}

%%K t12-19-1 p
19. --- Ik werd plotseling gewekt door de stem van een beambte die riep: « Deutsch-Avricourt!
\textsuperscript{(*)}. Iedereen uitstappen! » Er vond de inspectie van de douane plaats en
alle vervelende formaliteiten van de grens. Ik moest mijn zakhorloge gelijkzetten met de
\VUgras{grote klok} \textsuperscript{(89)} van het station, die vijfenvijftig minuten
voorliep; nauwelijks wakker, onderscheidde ik moeilijk de \VUgras{grote wijzer}
\textsuperscript{(90)} van de \VUgras{kleine} \textsuperscript{(91)}; de
\VUgras{wijzerplaat} \textsuperscript{(92)} zelf was helemaal verward door de
ochtendnevel.

%%K t12-20-1 p
20. --- Terwijl de douaniers hun inspectie deden, ontcijferde ik geeuwend de
\VUgras{aanplakbiljetten} \textsuperscript{(93)} die de stationsmuren versieren.
Ik ontbeet om de tijd te doden, raadpleegde de kaart van het Duitse
\VUgras{net} \textsuperscript{(94)} om te zien hoeveel afstand mij nog van
Straatsburg scheidt, en ik ging mijn wagen weer in, gesterkt door de hoop weldra het doel van
mijn reis te bereiken.
\VUsaut{6.0mm}
\VUfilet{22mm}
